\documentclass[11pt]{article}
\usepackage[margin=1in]{geometry}
\usepackage{amsmath,amssymb}
\usepackage{booktabs}
\usepackage[colorlinks=true,linkcolor=blue,citecolor=blue,urlcolor=blue]{hyperref}

\newcommand{\Wslow}{W_{\text{slow}}}
\newcommand{\Wfast}{W_{\text{fast}}}

\title{\textbf{BDH-Spike: A Native Spiking Neural Network Implementation of\\
the Baby Dragon Hatchling Architecture for\\
Neuromorphic Computing with Continual Learning}}
\author{Krzysztof Pika}
\date{Preprint, August 2026}

\begin{document}
\maketitle

\begin{abstract}
The Baby Dragon Hatchling (BDH) architecture replaces attention-softmax
machinery with brain-inspired neuron--synapse dynamics, yet its canonical
form still operates on continuous floating-point states. We present
\textbf{BDH-Spike}, a native Spiking Neural Network (SNN) realization of BDH
in which every activation inside the core layers is a discrete spike
$S(t)\in\{0,1\}$. The framework introduces (i) a Parametric Leaky
Integrate-and-Fire neuron with lateral BDH coupling (BDH-PLIF) trained by
surrogate gradients, (ii) a softmax-free, multiply-free attention mechanism
realized as associative AND-masking on binary spikes with an exact integer
bitwise inference path, (iii) a dual-weight plasticity engine separating
structural weights ($\Wslow$, trained offline by BPTT) from episodic weights
($\Wfast$, updated online by local three-factor Hebbian STDP under
\texttt{no\_grad}), and (iv) gradient-free homeostatic threshold regulation.
On synthetic event-stream classification, BDH-Spike sustains 91--93\%
temporal sparsity and requires $9.1\times$ fewer synaptic operations than
its dense equivalent. On a five-task split benchmark, episodic $\Wfast$
plasticity reduces measured catastrophic forgetting by ${\sim}24\%$ relative
to BPTT-only training. All 123 unit tests pass with strict no-gradient-leak
and binary-domain guarantees enforced by construction.
\end{abstract}

\noindent\textbf{Keywords:} spiking neural networks; neuromorphic computing;
Baby Dragon Hatchling; continual learning; STDP; surrogate gradients;
homeostasis; energy-efficient inference

\section{Introduction}
Large sequence models built on transformer machinery dominate machine
learning, but their energy profile is fundamentally hostile to edge and
neuromorphic hardware: dense multiply--accumulate (MAC) streams over
floating-point activations. The Baby Dragon Hatchling (BDH) architecture
replaces attention softmax with brain-inspired neuron--synapse dynamics,
yet its canonical formulation remains continuous-valued.

We argue that the natural substrate for BDH is the \emph{spiking} regime:
neurons integrate evidence over time, communicate through binary events, and
update synapses through local plasticity rules rather than global error
backpropagation. \textbf{BDH-Spike} operationalizes this claim as a complete,
tested software stack: a PLIF core with BDH coupling, softmax-free spiking
attention, a dual-weight plasticity engine, homeostatic regulation, full
model backbones, spike encoders, energy accounting, benchmarks, and
diagnostics.

Four architectural invariants govern every module: (1)~the \emph{binary
spike domain}: all activations in core layers satisfy $S(t)\in\{0,1\}$;
(2)~\emph{dual-weight plasticity}: structural weights $\Wslow$ are trained
offline via BPTT while episodic weights $\Wfast$ adapt online through local
STDP with zero global backpropagation; (3)~a \emph{canonical temporal
layout} $[T,B,C,\dots]$ preserved everywhere; and (4)~\emph{energy sparsity
first}, targeting 85--95\% silent states.

\section{The BDH-PLIF Neuron}
The core neuron extends Parametric Leaky Integrate-and-Fire with a lateral
BDH recurrent coupling vector $\mathbf{m}_{\text{BDH}}$ that nonlinearly
modulates excitability across time-steps:
\begin{align}
V[t] &= \beta\, V[t-1] + I_{\text{syn}}[t] + m_{\text{BDH}}[t-1]
        - S[t-1]\,V_{\text{th}}, \label{eq:mem}\\
S[t] &= \Theta\!\big(V[t] - V_{\text{th}}\big), \label{eq:spike}\\
m_{\text{BDH}}[t] &= \tanh\!\big(\rho\, m_{\text{BDH}}[t-1]
        + g\, S[t]\big), \label{eq:mbdh}
\end{align}
where $\beta=\sigma(\beta_{\text{logit}})$ is a per-channel learnable decay
constrained to $(0,1)$, $\rho$ the coupling decay and $g$ the coupling
strength. A hard reset forces $V[t]\leftarrow 0$ wherever a spike is
emitted. Differentiability is provided exclusively by the fast-sigmoid
surrogate applied at the threshold crossing,
\begin{equation}
\sigma'(x) \;=\; \frac{1}{\left(1 + k\,|x - V_{\text{th}}|\right)^2},
\qquad k = 25 .
\label{eq:surrogate}
\end{equation}
Forward values remain strictly binary; floating-point precision exists only
to carry gradients during offline training.

\section{Softmax-Free Spike-Driven Attention}
Attention is realized without Softmax, without scaling, and without
multiplies. Binary queries, keys and values are produced by thresholding
linear projections; association is counted by exact integer accumulation of
overlapping spikes; the mask is re-binarized and applied to values:
\begin{equation}
Q_s = \Theta(X W_q - V_{\text{th}}),\quad
K_s = \Theta(X W_k - V_{\text{th}}),\quad
V_s = \Theta(X W_v - V_{\text{th}}),
\end{equation}
\begin{equation}
C[n,m] = \sum_c Q_s[n,c] \wedge K_s[m,c], \qquad
A = \Theta(C - \theta_a), \qquad
Y = \Theta(A V_s - V_{\text{out}}).
\end{equation}
Two execution modes share one code path semantics: a \emph{surrogate} mode
(binary-valued floats carrying surrogate gradients into $W_q,W_k,W_v$) and
a \emph{bitwise} mode producing a strictly boolean/integer graph --- boolean
AND-mask, exact uint8/int16 accumulation --- mapping one-to-one onto
POPCOUNT/AC neuromorphic instructions. No float tensor is materialized after
input encoding in bitwise mode.

\section{Dual-Weight Plasticity}
Each synapse layer carries two weight matrices with disjoint learning
regimes. The fused transmission kernel is
\begin{equation}
I_{\text{syn}} = (\Wslow + \Wfast)^{\!\top} S_{\text{in}} .
\end{equation}
$\Wslow$ is an ordinary autograd parameter trained offline by BPTT through
Eq.~\eqref{eq:surrogate}. $\Wfast$ is a registered buffer updated
exclusively inside a \texttt{no\_grad} region by local three-factor Hebbian
(STDP) learning, following the canonical exponential trace recursion
\begin{align}
A_{\text{pre}}[t] &= A_{\text{pre}}[t-1]\,e^{-\Delta t/\tau_{\text{pre}}}
   + S_{\text{pre}}[t], \\
A_{\text{post}}[t] &= A_{\text{post}}[t-1]\,e^{-\Delta t/\tau_{\text{post}}}
   + S_{\text{post}}[t],
\end{align}
\begin{equation}
\Delta W[t] = M\cdot\Big(\eta_+\, S_{\text{post}}[t] \otimes A_{\text{pre}}[t]
   \;-\; \eta_-\, S_{\text{pre}}[t] \otimes A_{\text{post}}[t]\Big),
\label{eq:stdp}
\end{equation}
where $M$ is the third (modulatory) factor gating the write into $\Wfast$,
which is clamped to $|{\Wfast}|\le w_{\max}$ at every step. Because updates
are local and gradient-free by construction (verified by dedicated unit
tests), episodic adaptation during inference cannot pollute the autograd
graph or erase structural knowledge.

\paragraph{Homeostasis.} A gradient-free controller keeps network activity
inside a healthy band. With a per-layer firing-rate estimate
$r[t]=\beta_r r[t-1] + (1-\beta_r)\,\overline{S}[t]$,
\begin{equation}
V_{\text{th}}[t] = \mathrm{clip}\big(V_{\text{th}}[t-1]
   + \eta_v (r[t] - r_{\text{target}}),\; V_{\min}, V_{\max}\big).
\end{equation}

\section{Models and Encoders}
The stack assembles two backbones. \textbf{BDH-Spike-ViT} pipelines images
through a convolutional patch projection, min--max normalized steady-rate
encoding through the surrogate junction, softmax-free spiking attention, a
BDH-PLIF temporal readout over folded tokens, and mean-rate pooling into a
linear classifier. \textbf{BDH-Spike-Seq} is a streaming block fusing
dual-weight synapses with the PLIF cell, supporting online STDP and
homeostasis for continuous data. Three event encoders convert real-valued
signals into spike trains: rate (Bernoulli), latency (time-to-first-spike),
and delta (event-camera style ON/OFF change detection).

\section{Energy Accounting: SOPs vs.\ FLOPs}
Synaptic operations scale with active spikes rather than tensor density:
\begin{equation}
\text{SOPs} = \sum_{t=1}^{T} \mathrm{nnz}(S_{\text{in}}[t]) \times
\text{FanOut},
\end{equation}
against a dense equivalent of $2\,T B\,\text{FanIn}\,\text{FanOut}$ FLOPs.
Table~\ref{tab:energy} contrasts the regimes; Table~\ref{tab:measured}
reports measured results.

\begin{table}[h]
\centering
\caption{Dense ANN vs.\ BDH-Spike inference profile.}
\label{tab:energy}
\begin{small}
\begin{tabular}{lll}
\toprule
 & Dense ANN & BDH-Spike \\
\midrule
Core operation & MAC (32-bit float) & AC / POPCNT (bitwise) \\
Activation cost & every neuron, every step & emitting neurons only \\
Memory traffic & full weight matrix / step & active rows only \\
Inference-time learning & none & online STDP ($\Wfast$) \\
Target hardware & GPU / TPU & neuromorphic / edge MCU \\
\bottomrule
\end{tabular}
\end{small}
\end{table}

\section{Experiments}
\subsection{Event-Stream Classification}
On synthetic class-prototyped event streams ($T{=}8$, 5 classes,
batch~64), three epochs of surrogate-gradient training yield the energy
profile of Table~\ref{tab:measured}. Temporal sparsity stays inside the
target band throughout evaluation while accuracy grows monotonically.

\begin{table}[h]
\centering
\caption{Measured energy metrics (synthetic event benchmark, epoch 3).}
\label{tab:measured}
\begin{small}
\begin{tabular}{lr}
\toprule
Quantity & Measured value \\
\midrule
Temporal spike sparsity & 91--93\% \\
SOPs (evaluation pass) & $5.52\times10^{5}$ \\
Dense-equivalent FLOPs & $5.02\times10^{6}$ \\
FLOPs/SOP ratio & $\mathbf{9.1\times}$ \\
\bottomrule
\end{tabular}
\end{small}
\end{table}

\subsection{Split-Task Continual Learning}
Five sequential tasks share one hidden layer --- the standard condition for
catastrophic forgetting. Both schedules train the hidden features
identically by SGD; the dual-weight schedule additionally adapts $\Wfast$
online via Eq.~\eqref{eq:stdp} while each task stream flows through at
inference time. Forgetting is measured as the mean peak-to-final accuracy
drop over all tasks except the last:
\begin{equation}
F = \frac{1}{K-1}\sum_{k=1}^{K-1}
\Big(\max_j \mathrm{acc}_k[j] - \mathrm{acc}_k[K]\Big).
\end{equation}
Dual-weight plasticity reduces forgetting from $0.301$ to $0.230$
(${\sim}24\%$ relative reduction) under identical structural training.

\section{Software Guarantees}
The implementation enforces its invariants mechanically: 123 unit tests
assert binary-domain outputs, exact $[T,B,C]$ layouts, hard-reset semantics,
surrogate correctness against the closed form of Eq.~\eqref{eq:surrogate},
$\Wfast$ exclusion from \texttt{parameters()} with no gradient leak under
enabled grad mode, trace decay matching $e^{-\Delta t/\tau}$ exactly, and
weight clamping saturation. All tests pass in under six seconds on CPU.

\section{Conclusion}
BDH-Spike demonstrates that the Baby Dragon Hatchling architecture can be
realized natively in the spiking regime without sacrificing trainability or
continual adaptability: surrogate gradients preserve BPTT trainability, a
bitwise attention path targets neuromorphic silicon directly, and dual-weight
plasticity plus homeostasis deliver measurable continual-learning benefits
at 91--93\% activity sparsity. Future work includes Loihi~2/Lava export,
on-chip POPCOUNT kernels, and large-scale event-vision benchmarks.

\begin{thebibliography}{9}
\bibitem{bdh} Baby Dragon Hatchling (BDH): brain-inspired neuron--synapse
sequence-model architecture. arXiv preprint, 2025.
\bibitem{snntorch} J.~K. Eshraghian et al.
``Training Spiking Neural Networks Using Lessons From Deep Learning,''
\emph{Proceedings of the IEEE}, 2021.
\bibitem{spikingjelly} W.~Fang et al.
``SpikingJelly,'' \emph{Science Advances}, 2023.
\bibitem{loihi} M.~Davies et al.
``Loihi: A Neuromorphic Manycore Processor with On-Chip Learning,''
\emph{Proceedings of the IEEE}, 2018.
\bibitem{stdp} G.-Q. Bi and M.-M. Poo.
``Synaptic Modifications in Cultured Hippocampal Neurons,''
\emph{Journal of Neuroscience}, 1998.
\bibitem{surrogate} E.~O. Neftci, H.~Mostafa, and F.~Zenke.
``Surrogate Gradient Learning in Spiking Neural Networks,''
\emph{IEEE Signal Processing Magazine}, 2019.
\end{thebibliography}

\end{document}
